\originalchapter{18.901 官方作业}\label{ch:official-topology-assignment}

本章依据 MIT OCW 的 18.901 \emph{Introduction to Topology}, Fall 2004, James Munkres 课程作业页及其 Problem Set 5 原件整理. PS5 原件共两页: 第一页说明作业规则, 第二页是一张九行、十四列的空表. 前六行是具体空间或空间类别, 后三行是操作; 总计 $9\times14=126$ 个判定格. 以下中文题面依据原件翻译, 解答均为编者补充, 不是 MIT 发布的官方答案. 集合论采用 ZFC; 特别是任意积及基数比较按含选择公理的通常约定理解.

\originalsection{Problem Set 5 题面}
\begin{exercise}\label{ex:official-ps5}
本作业为选做, 用来补救任何一次成绩低于 $70$ 分的必做作业. 从下表选一行或多行, 在所选行的每个格内填 $+$ 表示“是”, 或填 $-$ 表示“否”. 每完整答对一行, 教师在当前最低分的那份作业上加 $6$ 分, 该份作业最多加至 $70$ 分. 同往常一样, 不允许合作. 可以在期末考试之前或考试时提交; 评判结果为最终结果. 原件以“No box tops needed to enter!”的玩笑结束这段说明, 意为不必凭商品包装上的兑奖标志参加.

前六行问:“给定空间是否满足该性质?” 后三行问:“给定操作是否保持该性质?” 为使原件中的课程记号可以独立阅读, 六行空间及三行操作说明如下.
\begin{enumerate}
\item $S_\Omega$: 所有可数序数构成的集合 $[0,\omega_1)$, 配以序拓扑. $\Omega=\omega_1$ 是第一个不可数序数, 不属于这个空间; $[0,\omega_1]$ 则是另一个包含最大元的空间.
\item 有序方形: 集合 $[0,1]\times[0,1]$ 按字典序排列, 即先比较第一坐标, 第一坐标相同才比较第二坐标, 再取这个线性序的序拓扑. 它与通常的欧氏方形拓扑不同.
\item $\R_\ell\times\R_\ell$: 每个 $\R_\ell$ 的下限拓扑以 $[a,b)$ ($a<b$) 为基, 两因子取积拓扑. 此空间通常称为 Sorgenfrey 平面.
\item $\R^\omega$ 的一致拓扑: $\omega$ 表示可数坐标集, 一致度量为 $\rho(x,y)=\sup_n\min\{1,|x_n-y_n|\}$. 不限制各序列本身有界.
\item $\R^I$ 的积拓扑: $I$ 是指标集, 基本开集只对有限个坐标施加开区间条件. 原件未限定 $|I|$, 所以答案必须区分空集、有限集、可数无限集及不可数集. 此行不是箱拓扑: 箱拓扑允许同时限制所有坐标. 第\ref{ch:products}章已说明二者的差别.
\item 任意度量空间: 问某性质是否由“是度量空间”这一条件保证, 并非挑选一个特定度量空间.
\item 取可数积: 各因子都有该性质时, 它们的可数积是否仍有该性质? 这里取积拓扑.
\item 取闭子空间: 空间有该性质时, 每个闭子空间是否都有该性质?
\item 取开子空间: 空间有该性质时, 每个开子空间是否都有该性质?
\end{enumerate}

十四列依原件顺序为以下性质. 原表没有零维或有限维列.
\begin{enumerate}
\item 连通: 不存在两个不交非空开集将空间分成两部分.
\item 道路连通: 任意两点之间都有连续道路 $\gamma:[0,1]\to X$.
\item 紧致: 每个开覆盖都有有限子覆盖.
\item 局部紧 Hausdorff: 空间 Hausdorff, 且每点有一个紧邻域. 在 Hausdorff 情形, 这等价于每点有开邻域, 其闭包紧.
\item Hausdorff: 两个不同点有互不相交的开邻域. 这是原件中独立的一列.
\item 正则: 按第\ref{ch:separation}章约定包含 $T_1$, 且点与不含它的闭集可以用不交开邻域分离.
\item 正规: 包含 $T_1$, 且不交闭集可以用不交开邻域分离.
\item 第一可数: 每点有可数邻域基.
\item 第二可数: 有可数开集基.
\item Lindel\"of: 每个开覆盖都有可数子覆盖.
\item 有可数稠密子集: 即可分.
\item 局部可度量化: 每点有一个开邻域, 该邻域的子空间拓扑可由某个度量生成. 此列不另加全空间 Hausdorff 条件.
\item 可度量化: 全空间拓扑可由某个度量生成.
\item 完全正则: 包含 $T_1$, 且对每个闭集 $F$ 和 $x\notin F$, 存在连续函数 $f:X\to[0,1]$ 满足 $f(x)=1$ 和 $f|_F=0$.
\end{enumerate}
后三行的肯定项是对所有满足所讨论性质的空间作出的断言; 不为其他列额外假定 Hausdorff、正则或非空. 空空间视为连通、道路连通且紧致, 这些约定不改变表中的反例.
\end{exercise}

\originalsection{判定与证明}
\begin{solution}
下面两张表按原顺序各列七项, 合起来覆盖全部 $126$ 格. 对 $\R^I$ 行, $\mathrm E$ 表示“当且仅当 $I=\varnothing$”, $\mathrm F$ 表示“当且仅当 $I$ 有限”, $\mathrm C$ 表示“当且仅当 $I$ 至多可数”, $\mathrm D$ 表示“当且仅当 $|I|\le\mathfrak c$”, 其中 $\mathfrak c=|\R|$. 这些条件是相应格中 $+$ 的准确条件, 条件不成立即填 $-$. 其他行的 $+$ 与 $-$ 均无附加条件.

\begin{center}
\setlength{\tabcolsep}{4pt}
\begin{tabular}{lccccccc}
\hline
空间或操作 & 连通 & \shortstack{道路\\连通} & 紧致 & \shortstack{局部紧\\Hausdorff} & Hausdorff & 正则 & 正规\\
\hline
$S_\Omega$ & $-$ & $-$ & $-$ & $+$ & $+$ & $+$ & $+$\\
有序方形 & $+$ & $-$ & $+$ & $+$ & $+$ & $+$ & $+$\\
$\R_\ell\times\R_\ell$ & $-$ & $-$ & $-$ & $-$ & $+$ & $+$ & $-$\\
$\R^\omega$ 一致拓扑 & $-$ & $-$ & $-$ & $-$ & $+$ & $+$ & $+$\\
$\R^I$ 积拓扑 & $+$ & $+$ & $\mathrm E$ & $\mathrm F$ & $+$ & $+$ & $\mathrm C$\\
任意度量空间 & $-$ & $-$ & $-$ & $-$ & $+$ & $+$ & $+$\\
取可数积 & $+$ & $+$ & $+$ & $-$ & $+$ & $+$ & $-$\\
取闭子空间 & $-$ & $-$ & $+$ & $+$ & $+$ & $+$ & $+$\\
取开子空间 & $-$ & $-$ & $-$ & $+$ & $+$ & $+$ & $-$\\
\hline
\end{tabular}
\end{center}

\begin{center}
\setlength{\tabcolsep}{4pt}
\begin{tabular}{lccccccc}
\hline
空间或操作 & \shortstack{第一\\可数} & \shortstack{第二\\可数} & Lindel\"of & 可分 & \shortstack{局部\\可度量} & 可度量 & \shortstack{完全\\正则}\\
\hline
$S_\Omega$ & $+$ & $-$ & $-$ & $-$ & $+$ & $-$ & $+$\\
有序方形 & $+$ & $-$ & $+$ & $-$ & $-$ & $-$ & $+$\\
$\R_\ell\times\R_\ell$ & $+$ & $-$ & $-$ & $+$ & $-$ & $-$ & $+$\\
$\R^\omega$ 一致拓扑 & $+$ & $-$ & $-$ & $-$ & $+$ & $+$ & $+$\\
$\R^I$ 积拓扑 & $\mathrm C$ & $\mathrm C$ & $\mathrm C$ & $\mathrm D$ & $\mathrm C$ & $\mathrm C$ & $+$\\
任意度量空间 & $+$ & $-$ & $-$ & $-$ & $+$ & $+$ & $+$\\
取可数积 & $+$ & $+$ & $-$ & $+$ & $-$ & $+$ & $+$\\
取闭子空间 & $+$ & $+$ & $+$ & $-$ & $+$ & $+$ & $+$\\
取开子空间 & $+$ & $+$ & $-$ & $+$ & $+$ & $+$ & $+$\\
\hline
\end{tabular}
\end{center}

先说明会重复使用的判据. 命题\ref{thm:second-countable}给出第二可数蕴含第一可数、可分和 Lindel\"of; 定理\ref{thm:metric-countability}给出度量空间中第二可数、可分与 Lindel\"of 等价. 命题\ref{thm:normal-examples}给出度量空间正规、紧 Hausdorff 空间正规, 以及正则 Lindel\"of 空间正规. 正规空间由 Urysohn 引理\ref{thm:urysohn}完全正则. Hausdorff 空间中紧集闭, 见引理\ref{thm:compact-hausdorff}.

完全正则性传给子空间及任意积的细节也在这里固定. 若 $A\subseteq X$, $F\subseteq A$ 相对闭, 将 $F$ 写成 $A\cap H$, 其中 $H$ 在 $X$ 中闭; 对 $x\in A\setminus F$ 用 $X$ 上分离 $x,H$ 的函数再限制到 $A$. 积空间中, 对 $x\notin F$ 取基本矩形 $\bigcap_{i\in J}\pi_i^{-1}(U_i)$ 包含 $x$ 并避开 $F$, 其中 $J$ 有限. 在每个因子取 $f_i(x_i)=1$, $f_i=0$ 于 $X_i\setminus U_i$. 函数 $\prod_{i\in J}f_i\circ\pi_i$ 连续, 在 $x$ 处为 $1$, 在 $F$ 上为 $0$. $T_1$ 性同样传给子空间与积. 完全正则蕴含正则: 上述函数的开集 $\{f>2/3\}$ 与 $\{f<1/3\}$ 分别包含该点和该闭集并互不相交. 下述证明依次对应原件九行.

\begin{enumerate}
\item \textbf{$S_\Omega$.}
每个初始区间 $[0,\alpha]$ 都紧. 事实上, 对它的任一开覆盖, 用超限归纳证明每个 $[0,\beta]$ 都有有限子覆盖: 后继步骤加一个覆盖后继点的开集; 极限步骤中, 覆盖 $\beta$ 的开集包含某个 $(\gamma,\beta]$, 再覆盖 $[0,\gamma]$ 即可. 这也证明 $[0,\omega_1]$ 紧. 序拓扑 Hausdorff; 每个 $\alpha$ 的开邻域 $[0,\alpha+1)$ 就是紧开集 $[0,\alpha]$, 所以局部紧 Hausdorff.

对每个 $\alpha<\omega_1$, $[0,\alpha]$ 可数, 其序拓扑有可数基, 并且紧 Hausdorff 而正则; 定理\ref{thm:metrization}使它可度量化. 上述开邻域证明局部可度量化, 从而第一可数. 集合 $[0,\alpha]$ 及其补集都开, 可用于分离任意两个不同序数, 所以空间不连通, 且任意连通子集至多一个点; 因道路像连通, 它也不道路连通.

开覆盖 $\{[0,\alpha):\alpha<\omega_1\}$ 没有可数子覆盖: 可数个可数序数的上确界仍小于 $\omega_1$. 同理每个可数点集有界而不稠密. 所以不紧、非 Lindel\"of、不可分且非第二可数.

初始区间 $[0,\alpha]$ 闭开, 其差集 $(\beta,\alpha]=[0,\alpha]\setminus[0,\beta]$, 连同含最小元的初始区间, 形成闭开邻域基, 从而得到正则性. 若两闭集 $A,B$ 都无界, 可交替选择严格递增的 $a_n\in A,b_n\in B$, 其可数上确界 $\lambda<\omega_1$ 同时是两条子序列的极限; 闭性给出 $\lambda\in A\cap B$. 因此不交闭集至少有一个有界. 有界闭集位于某个紧初始区间中, 因而紧. 正则性分离这个紧集的每一点与另一个闭集, 取有限个邻域并, 即将两闭集分离. 故正规并完全正则.

最后, $S_\Omega$ 可数紧: 任一可数开覆盖若无有限子覆盖, 可选 $x_n$ 避开前 $n$ 个成员; 所有 $x_n$ 落在某个紧 $[0,\beta]$ 中, 此区间上的有限子覆盖与选择矛盾. 可数紧的度量空间紧: 每个序列都有聚点, 否则取一个无聚点的无限值域, 其补集及分别只含一个值域点的开邻域组成没有有限子覆盖的可数开覆盖; 第一可数性再给收敛子序列, 用第\ref{ch:compact}章的紧性等价. 因本空间不紧, 它不可能可度量化.

\item \textbf{有序方形.}
字典序是稠密序: 同一竖直纤维内可取第二坐标的中间值; 不同纤维之间可取第一坐标的中间值. 它还有最小元 $(0,0)$、最大元 $(1,1)$, 每个非空子集有上确界. 具体地, 先取第一坐标的上确界 $a$; 若集合在纤维 $a$ 中有点, 再取那些点第二坐标的上确界, 否则取 $(a,0)$.

完备且有两端点的序拓扑紧. 对开覆盖令 $T$ 为所有使初始区间至该点有有限子覆盖的点; 覆盖最小元的成员保证 $T$ 非空. 取 $c=\sup T$. 覆盖 $c$ 的开邻域含其左侧一段; 与 $T$ 中足够靠近 $c$ 的点的有限子覆盖合并, 即覆盖至 $c$, 所以 $c\in T$. 若 $c$ 不是最大元, 同一开邻域还覆盖右侧一段, 与 $c=\sup T$ 矛盾. 因而全空间有有限子覆盖. 序拓扑 Hausdorff, 故空间紧、局部紧 Hausdorff、正规、正则、Lindel\"of 和完全正则.

稠密完备序的区间连通: 若 $x<y$ 分别属于不交开集构成的分离两侧, 考察左侧在 $[x,y]$ 内的上确界 $c$. 若 $c$ 在左侧, 其开邻域和序稠密性给出更大的左侧点; 若在右侧, 其左邻域排除逼近 $c$ 的左侧点. 端点情形因 $x,y$ 分居两侧也排除. 所以有序方形连通.

集合 $\{x\}\times(0,1)$ ($x\in[0,1]$) 是不可数族两两不交非空开集, 可数稠密集不可能逐一与之相交. 故不可分、非第二可数且不可度量化, 最后一项也可由紧度量空间第二可数推出. 每点仍有可数邻域基: 纤维内部只用第二坐标的可数区间; 在 $(x,0)$ 或 $(x,1)$ 处, 从相邻纤维取第一坐标逐步逼近 $x$ 的序列, 再配合同一纤维内的第二坐标序列即可. 两个全局端点只取存在的那一侧.

若一道路的端点第一坐标不同, 其连通像必须包含两端点之间的整个序区间, 因而遇到不可数多个竖直开区间. 但道路像可分, 因为 $[0,1]$ 中有理点的像稠密, 与这些不交开集矛盾. 因此不道路连通. 在 $(1/2,0)$ 的任意开邻域内, 都能找到一个包含不可数多个完整纤维的闭序区间. 此区间紧而不可分; 若该邻域可度量化, 此区间将是紧度量空间并可分, 矛盾. 故空间也不局部可度量化.

\item \textbf{Sorgenfrey 平面.}
每个 $[a,b)$ 在 $\R_\ell$ 中既开又闭, 基本矩形也闭开. 因而平面 Hausdorff、正则, 且完全正则: 对点与闭集选择避开闭集的闭开矩形, 用其示性函数分离. 闭开矩形分离不同点, 所以任意连通子集至多一点, 不连通也不道路连通. 点 $(x,y)$ 的 $[x,x+1/n)\times[y,y+1/n)$ 给出第一可数性. $\mathbb Q^2$ 稠密.

反对角线 $A=\{(t,-t):t\in\R\}$ 是闭离散子空间. 在其点 $(t,-t)$, 一个下限基本矩形只与 $A$ 相交于该点. 对不在 $A$ 上的点, 按坐标和为正或负取足够小的矩形避开 $A$, 故 $A$ 闭; 其每个子集也在平面中闭. 如果平面正规, Urysohn 引理会为每个 $B\subseteq A$ 提供一个连续 $[0,1]$ 值函数, 在 $B$ 上为 $0$, 在 $A\setminus B$ 上为 $1$. 这些函数互不相同, 至少有 $2^{\mathfrak c}$ 个. 然而连续实值函数由其在可数稠密集上的值唯一决定, 最多只有 $\mathfrak c^{\aleph_0}=\mathfrak c$ 个, 与 Cantor 定理矛盾. 因而平面不正规. 正则 Lindel\"of 空间正规, 故它非 Lindel\"of; 第二可数蕴含 Lindel\"of, 故也非第二可数. 由非正规还知不紧和不可度量化.

它不局部紧. 若某点有紧邻域 $K$, 则第一投影 $C=\pi_1(K)$ 是下限直线的紧集, 并包含一个非空通常开区间. 恒等映射 $C\to C$ 的通常实数子空间拓扑是连续双射, 紧到 Hausdorff 的判据使它为同胚. 但在该区间内部取 $c<d$, 下限开集 $[c,d)\cap C$ 不是通常相对开集, 矛盾.

它不局部可度量化. 任意非空基本矩形都包含一条相对闭的水平线段, 同胚于下限区间 $[a,b)$. 这个区间可分却非第二可数: 对每个 $t\in[a,b)$, 在一个假定的基中选择 $t\in B_t\subseteq[t,b)$, 不同 $t$ 所选 $B_t$ 不同. 由度量空间的可数性等价, 这个区间不能可度量化; 所以任何包含矩形的开邻域也不能可度量化.

\item \textbf{$\R^\omega$ 的一致拓扑.}
$\rho$ 是度量, 所以 Hausdorff、正规、正则、完全正则、第一可数且局部可度量化. 对每个 $x$, 集合 $x+\ell^\infty$ 既开又闭: 小于 $1$ 的一致球中坐标差一致有界, 不同的有界差等价类互不相交并构成开分划. 例如零序列与 $(n)_n$ 不在同一类. 因而空间不连通, 不道路连通. 同一类内的线性道路确实连续, 因为其坐标差有统一有限界.

$\{0,1\}^{\mathbb N}$ 中不同点距离为 $1$, 是不可数分离集. 半径 $1/3$ 的球两两不交, 所以没有可数稠密集; 由度量可数性等价也非第二可数且非 Lindel\"of. 它不局部紧: 任何点的任何紧邻域都包含某个球 $B_\rho(x,r)$, 取 $0<t<\min\{r,1\}$, 则其中的 $x+t e_n$ 两两距离为 $t$, 没有收敛子序列, 与紧度量空间的序列紧性矛盾. 因而也不紧.

\item \textbf{$\R^I$ 的积拓扑.}
任意两点的逐坐标线性道路 $\gamma(t)_i=(1-t)x_i+ty_i$ 连续, 因为每个坐标连续; 所以道路连通且连通. 完全正则性由前述乘积函数证明, 也得到正则和 Hausdorff. 若 $I=\varnothing$, 积是一个点, 全部十四项都成立. 若 $I\ne\varnothing$, 投影到一个 $\R$ 因子连续且满射, 所以不紧.

有限 $I$ 的积是欧氏空间, 局部紧 Hausdorff. 无限 $I$ 时, 每个基本邻域都留下一个不受限制的坐标. 若有紧邻域 $K$, 它包含一个基本邻域, 在这个坐标的投影就是 $\R$, 而紧集投影不能为 $\R$, 矛盾. 因而局部紧性恰好对应有限 $I$.

至多可数 $I$ 时, 命题\ref{thm:countable-product-metric}给出可度量化和第二可数, 于是第一可数、正规、Lindel\"of 且局部可度量化. 不可数 $I$ 时, 任一点都不第一可数: 假设有可数邻域基, 将其各缩成一个基本邻域, 全部受约束坐标的并仍可数; 取未受约束的 $j$, 这些邻域不能包含于 $\{y:|y_j-x_j|<1\}$. 所以非第二可数、不可度量化且不局部可度量化. 后一项使用局部可度量化必推出第一可数.

可分性的准确界是 $|I|\le\mathfrak c$. 必要性: 对可数稠密集 $D$, 各坐标在 $D$ 上的取值向量必须两两不同, 否则 $D$ 落在两个坐标相等的闭集内而不稠密. 因此 $I$ 注入 $\R^D$, 后者至多有 $\mathfrak c$ 个元素. 充分性: 为各 $i$ 选择不同的实数 $t_i$, 取可数集 $\{(p(t_i))_{i\in I}:p\in\mathbb Q[t]\}$. 给定有限坐标的开区间条件, 用 Lagrange 插值使一个实系数多项式在那些不同的 $t_i$ 处取指定的区间内部值, 再充分逼近其有限个系数为有理数, 仍满足全部条件. 故这个可数集稠密.

不可数积的非正规性需要真正的反例, 不能仅从不可度量化推出. 下面完整给出 Stone 反例的有限坐标证明. 在 $X=\mathbb N^{\omega_1}$ 中取 $\mathbb N=\{0,1,2,\ldots\}$ 为离散空间, 记 $[F,x]=\{y:y|_F=x|_F\}$, 其中 $F$ 有限. 令 $H_j$ ($j=0,1$) 由那些除数值 $j$ 外每个数值至多出现一次的函数组成. 若一函数不在 $H_j$, 用两个坐标同取某个 $m\ne j$ 的柱集就能见证, 所以 $H_j$ 闭. 若函数同时属于 $H_0,H_1$, 每个数值都至多出现一次, 成为 $\omega_1\to\mathbb N$ 的单射, 不可能. 因此两闭集不交.

任取开集 $U\supseteq H_0$, $V\supseteq H_1$. 递归构造有限集 $F_n$ 和相容单射 $q_n:F_n\to\{2,3,\ldots\}$. 从 $F_0=\varnothing$ 开始, 令 $f_n$ 在 $F_n$ 上等于 $q_n$, 外面为 $0$. 因 $f_n\in H_0$, 可取有限 $F_{n+1}\supseteq F_n$ 使 $[F_{n+1},f_n]\subseteq U$, 再把 $q_n$ 扩成单射 $q_{n+1}$. 令 $F=\bigcup_nF_n$, $q=\bigcup_nq_n$, 并令 $g$ 在 $F$ 上等于 $q$, 外面为 $1$; 则 $g\in H_1$. 取有限 $G$ 使 $[G,g]\subseteq V$, 再取 $n$ 使 $G\cap F\subseteq F_n$. 在 $G\cap F_{n+1}$ 上, $g$ 与 $f_n$ 一致, 因为这个交集落在 $F_n$ 内. 两组有限坐标条件相容, 所以 $[G,g]\cap[F_{n+1},f_n]\ne\varnothing$. 得 $U\cap V\ne\varnothing$, 证明 $X$ 非正规.

对不可数 $I$, 在 ZFC 中选择 $J\subseteq I$, $|J|=\aleph_1$. 闭子空间 $\mathbb N^J\times\{0\}^{I\setminus J}$ 同胚于上述 $X$; 正规性传给闭子空间, 所以 $\R^I$ 非正规. 它已是正则空间, 正则 Lindel\"of 空间正规, 故它也非 Lindel\"of. 这同时完成两格的不可数情形, 无须连续统假设. Stone 反例的来源见本章末尾; 上面的有限坐标论证已给出所需全部步骤.

\item \textbf{任意度量空间.}
度量空间 Hausdorff、正规和正则, 由距离函数或 Urysohn 引理完全正则; 球 $B(x,1/n)$ 给第一可数性. 定义本身给可度量化及局部可度量化. 两点离散空间反驳普遍的连通性与道路连通性; $\R$ 反驳普遍紧性; 上一行的一致度量空间反驳普遍局部紧性. 不可数离散度量空间不第二可数、不可分且非 Lindel\"of: 稠密集必须是全集, 单点开覆盖没有可数子覆盖.

\item \textbf{取可数积.}
先说明有限积连通: 两因子积中, 固定 $(a,b)$, 各集合 $(X\times\{b\})\cup(\{x\}\times Y)$ 是两连通切片的并, 共同含 $(a,b)$; 对 $x$ 取并得到全积, 再归纳到有限积. 可数积的连通性由有限坐标改变点的集合证明. 固定积点 $a$, 只改变有限坐标的各有限子积都连通, 共同含 $a$, 故其并连通; 该并在积中稠密. 连通集的闭包连通, 因为闭包若有分离, 两个非空相对开侧都会遇到原稠密连通集并将其分离. 所以全积连通. 因子或积为空时结论直接成立. 道路连通性由逐坐标道路及积的连续性判据得到. 紧性由 Tychonoff 定理\ref{thm:tychonoff}得到; Hausdorff 性按不同坐标分离. 正则性则在有限受约束坐标逐一取闭包包含于给定坐标开集的小邻域, 其基本矩形闭包仍包含在原矩形中. 第一可数的局部基由有限坐标的可数局部基组合而成; 第二可数和可度量化见命题\ref{thm:countable-product-metric}. 可分性可取各因子的可数稠密集, 固定其各一点, 只在有限坐标改取稠密集中的点; 所得点集可数且稠密. 完全正则性已由有限个分离函数的乘积证明.

局部紧 Hausdorff 性不保持: $\R^{\mathbb N}$ 即为反例. 正规性和 Lindel\"of 性不保持, 均可用 $\R_\ell\times\R_\ell$; 若要求写成可数无限积, 补上单点因子即可. 为确认因子性质, 这里证明 $\R_\ell$ 是 Lindel\"of. 给定开覆盖, 为每个 $x$ 取 $[x,b_x)$ 包含于某个覆盖成员. 通常开集 $O=\bigcup_x(x,b_x)$ 可从这些通常开区间中选可数子覆盖, 因通常实线第二可数. 对 $x\in\R\setminus O$, 区间 $(x,b_x)$ 两两不交: 若 $x<y$ 均在此补集, 必有 $b_x\le y$. 各区间选一个有理数, 知补集可数. 用覆盖 $O$ 的可数个区间对应的原覆盖成员, 再加补集各点对应的成员, 得到下限直线的可数子覆盖. 下限直线正则, 因而正规.

局部可度量化也不保持: $S_\Omega^{\mathbb N}$ 的每个因子局部可度量化, 但任何非空基本邻域都留下一个完整 $S_\Omega$ 因子. 固定其他坐标, 它含一个同胚于 $S_\Omega$ 的子空间, 因而不能可度量化. 若积中某点有可度量开邻域, 将其缩成基本邻域就得到矛盾.

\item \textbf{取闭子空间.}
$\{0,1\}\subseteq\R$ 是闭子空间, 但不连通也不道路连通. 紧性由闭子集紧得到; 正规性因为闭子空间中的闭集仍在母空间中闭而保持. Lindel\"of 性由给相对开覆盖添上闭子空间的补集得到母空间开覆盖, 再取可数子覆盖. Hausdorff、正则、第一可数、第二可数、可度量化、局部可度量化和完全正则均传给任意子空间: 限制分离邻域、邻域基、基、度量或分离函数即可; 正则性的闭包收缩论证见第\ref{ch:separation}章练习. 闭子空间的局部紧 Hausdorff 性, 用原紧邻域与闭子空间相交得到.

可分性不保持: Sorgenfrey 平面可分, 但其闭离散反对角线 $A$ 不可数, 不是可分空间. 此反例说明不能简单把母空间的可数稠密集与闭子空间取交.

\item \textbf{取开子空间.}
$\R\setminus\{0\}$ 不连通也不道路连通; $(0,1)\subseteq[0,1]$ 说明紧性不保持. 第一可数、第二可数、Hausdorff、正则、可度量化、局部可度量化及完全正则按上一项传给所有子空间. 可分性在开子空间保持: 若 $D$ 可数稠密且 $U$ 开, 每个非空相对开集也是母空间的非空开集, 因而与 $D\cap U$ 相交. 局部紧 Hausdorff 也保持: 给定 $x\in U$, 用局部紧 Hausdorff 的邻域收缩引理\ref{lem:local-compact-shrink}取得闭包紧且包含于 $U$ 的小开邻域.

Lindel\"of 性不保持: $[0,\omega_1)$ 是紧空间 $[0,\omega_1]$ 的开子空间, 已知非 Lindel\"of. 正规性不保持的反例是删除角点的 Tychonoff 板. 令 $P=[0,\omega_1]\times[0,\omega]$, 两因子取序拓扑. 它们紧 Hausdorff, 所以 $P$ 紧 Hausdorff 并正规; $Y=P\setminus\{(\omega_1,\omega)\}$ 开. 在 $Y$ 中, 集合 $A=\{\omega_1\}\times[0,\omega)$ 和 $B=[0,\omega_1)\times\{\omega\}$ 不交且闭, 因其在 $P$ 中的闭包只额外添入已删除的角点. 若开集 $U\subseteq Y$ 包含 $A$, 对每个有限序数 $n$ 可取 $\alpha_n<\omega_1$ 使 $(\alpha_n,\omega_1]\times\{n\}\subseteq U$. 令 $\alpha=\sup_n\alpha_n<\omega_1$, 取 $\beta=\alpha+1$. 点 $(\beta,\omega)\in B$ 的每个邻域都含 $(\beta,n)$, 其中 $n$ 足够大; 而这些点全在 $U$. 因此任何包含 $B$ 的开集都与 $U$ 相交, $Y$ 不正规.
\end{enumerate}
\end{solution}

\originalsection{官方作业目录与题源}
\label{sec:official-topology-directory}
官方作业页将 weekly exercises 与 problem sets 区分开: weekly exercises 不交评分, 用于备考, 鼓励合作; PS0 是集合论语言的诊断作业, 其分数只供选课建议; PS1--4 要求独立提交规范书面解答; PS5 原件明确标为 optional. 下列编号均指 James R. Munkres, \emph{Topology}, 2nd ed., Prentice-Hall, 1999 的节号和习题号. 官方页对 PS0--4 及 weekly exercises 仅给教材题号, \textbf{缺少可公开下载的完整题面}; 本章只列目录, 不据题号重构或转录版权教材题面. PS5 则有完整两页原件, 已在习题\ref{ex:official-ps5}翻译并解答.

\begin{description}
\item[PS0.] \S1: 2, 5; \S2: 4(c), 4(e). 均只给答案. \S1 第2题的 (j), (k), (l) 答“是”或“否”, 其余答所列逻辑或包含关系符号. 原件未在作业页重刊这道教材题的各子问.
\item[PS1.] \S3: 13; \S9: 8; \S13: 7; \S17: 16, 18. \S13 第7题及 \S17 第16、18题只给答案.
\item[PS2.] \S18: 13; \S20: 6(a)(b), 8(b)(c); \S24: 4. \S20 第8题假定 (a), 其中 (c) 只给答案.
\item[PS3.] \S26: 9, 12; \S28: 4; \S30: 5.
\item[PS4.] \S31: 7(a)(b)(d); \S33: 4; \S34: 4, 5; \S38: 9. \S34 第4、5题给证明或反例.
\item[PS5.] 官方完整 PDF, 九行十四列, 选做.
\end{description}

\begin{description}
\item[第1周: 逻辑与基础.] \S1: 3; \S2: 1, 2, 4, 5.
\item[第2周: 关系、基数、选择公理.] \S3: 11, 12, 15; \S5: 4, 5; \S6: 3, 6; \S7: 3, 4, 5.
\item[第3周: 拓扑、闭集.] \S7: 6; \S13: 2, 6, 8; \S16: 3, 4, 8, 10; \S17: 3, 4, 5, 6, 8, 9.
\item[第4周: 连续函数、任意积.] \S17: 10, 11, 12; \S18: 2, 3, 7, 8; \S19: 1, 2, 3, 6, 8.
\item[第5周: 度量拓扑.] \S19: 7; \S20: 2, 4, 5, 8(a); \S21: 3, 4, 6, 7.
\item[第6周: 商拓扑.] \S22: 2, 3, 6.
\item[第7周: 连通空间、紧空间.] \S23: 2, 3, 5, 7, 8; \S24: 1, 5, 8, 9; \S25: 1, 2, 3; \S26: 1, 4, 5, 6.
\item[第8周: 紧性续讲.] \S26: 7, 8; \S27: 1, 4; \S28: 1, 2, 3.
\item[第9周: 良序集、极大原理; 期中考试.] 官方页未列习题号.
\item[第10周: 可数性与分离公理.] \S10: 2, 3, 6; \S29: 5, 6, 7, 8; \S30: \texttt{1, 2, 4 7, 8, 9}. 官方页面在 $4$ 和 $7$ 之间未印分隔符, 此处照录; 不将这一处排印歧义视为已核实的独立题号分隔.
\item[第11周: Urysohn 引理、可度量化.] \S31: 1, 3, 6; \S32: 1, 2, 3, 6, 7.
\item[第12周: Tietze 定理.] \S33: 2, 6, 7, 8.
\item[第13周: Tychonoff 定理、Stone--\v{C}ech 紧化.] \S34: 1, 3, 7, 8; \S37: 2, 3; \S38: 5, 6.
\item[第14周: Baire 空间、维数理论.] \S38: 3, 4; \S36: 1, 5; \S48: 1, 2, 3, 5, 6. 保留官方页面的节号顺序.
\item[第15周: 嵌入欧氏空间.] 官方页未列习题号. 随后的期末考试行也未列习题号.
\end{description}

官方目录与 PS5 下载链接见
\href{https://ocw.mit.edu/courses/18-901-introduction-to-topology-fall-2004/pages/assignments/}{MIT OCW: 18.901 Fall 2004 Assignments}.
关于不可数积的非正规性, 原始出处是 A. H. Stone, ``Paracompactness and product spaces'', \emph{Bulletin of the American Mathematical Society} \textbf{54} (1948), 977--982. 闭集 $H_0,H_1$ 的定义及这一反例的不同证明也见 N. Noble, \href{https://arxiv.org/abs/2002.02483}{\emph{Poorly Separated Infinite Normal Products}}, \S6.1. 本章已展开有限坐标证明, 并未将这一结论视为前十五章已经证明的教材结论.
